%% LyX 2.4.4 created this file.  For more info, see https://www.lyx.org/.
%% Do not edit unless you really know what you are doing.
\documentclass[english]{article}
\usepackage[T1]{fontenc}
\usepackage[latin9]{inputenc}
\usepackage{units}
\usepackage{enumitem}
\usepackage{amsmath}
\usepackage{amsthm}
\usepackage{amssymb}
\usepackage{geometry}
\geometry{verbose,tmargin=1in,bmargin=1in,lmargin=1in,rmargin=1in}

\makeatletter
%%%%%%%%%%%%%%%%%%%%%%%%%%%%%% Textclass specific LaTeX commands.
\numberwithin{equation}{section}
\numberwithin{figure}{section}
\newlength{\lyxlabelwidth}      % auxiliary length 

%%%%%%%%%%%%%%%%%%%%%%%%%%%%%% User specified LaTeX commands.
\usepackage{fancyhdr}
\usepackage{lastpage}
\pagestyle{fancy}
\usepackage{enumitem}
\usepackage{ifthen}
\everymath=\expandafter{\the\everymath\displaystyle}
%\DeclareMathSizes{10}{10}{10}{10}
\usepackage{multicol}

\usepackage{tikz}
\usetikzlibrary{patterns}
\usetikzlibrary{plotmarks}
\usepackage{pgfplots}
\pgfplotsset{compat=newest}

\usepackage{calc}
\usepackage[nomessages]{fp}

\usepackage{datetime}
% Added by lyx2lyx
\setlength{\parskip}{\smallskipamount}
\setlength{\parindent}{0pt}

\makeatother

\usepackage{babel}
\begin{document}
\renewcommand{\author}{Dr.\,James Ashe}

\newcommand{\classNum}{336}

\newcommand{\classSec}{A}

\newcommand{\examType}{HW 1.1 Solutions}

\renewcommand{\date}{\formatdate{26}{8}{2013}}

%%First page's header%%%%%%%%%%%%%%%%%%%%%%
\fancypagestyle{firststyle} %%header settings for first pagr
{
	\fancyhead[L]{MTH \classNum{}(\classSec{}) \author{}\\
	\textbf{\examType{}} ~\date{}}
	\fancyhead[C]{}
	\fancyhead[R]{\textbf{Name:\hspace{4cm}}}
	\setlength{\headsep}{14pt}
}
%%%%%%%%%%%%%%%%%%%%%%%%%%%%%%%%%%%%%%%%%%

%%Next page's header%%%%%%%%%%%%%%%%%%%%%%%
\setlist{nolistsep}
\lhead{\classNum{}(\classSec{})} 
\chead{\textbf{Name:}} 
\cfoot{\thepage{}~of~\pageref{LastPage}} 
\setlength{\headsep}{5pt} 
%%%%%%%%%%%%%%%%%%%%%%%%%%%%%%%%%%%%%%%%%%%

%%Various settings; see the preamble for other settings%%%%%
%\setenumerate[0]{leftmargin=12pt,itemindent=0pt,labelwidth=0pt}
\setlist{topsep=.5pt}
\renewcommand{\labelenumi}{\textbf{\arabic{enumi}.}}
\renewcommand{\labelenumii}{\textbf{\alph{enumii}.}}
\thispagestyle{firststyle}
%%%%%%%%%%%%%%%%%%%%%%%%%%%%%%%%%%%%%%%%%%

\newcommand{\xmax}{0}
\newcommand{\xmin}{-\xmax}
\newcommand{\xstep}{0}
\newcommand{\ymax}{0}
\newcommand{\ymin}{-\ymax}
\newcommand{\ystep}{0} 
\newcommand{\dspace}{\vspace{1cm}}
\newcommand{\xa}{0}
\newcommand{\xb}{0}
\newcommand{\ya}{1}
\newcommand{\yb}{1}
\begin{enumerate}[resume]
\item \textbf{$\mathbf{x=}\left(2,3\right)$} and \textbf{$\mathbf{y=}(-1,1)$}\begin{multicols}{3}

\begin{enumerate}[resume]
\item $\mathbf{x+y}=(1,4)$\\
\renewcommand{\ya}{1}
\renewcommand{\yb}{4}\renewcommand{\xmax}{5}
\renewcommand{\xmin}{-5}
\renewcommand{\xstep}{0} %must be xmin+n
\renewcommand{\ymax}{5}
\renewcommand{\ymin}{-5}
\renewcommand{\ystep}{0} %must be ymin+n
%%%%%%%
\begin{tikzpicture} 
\begin{axis}[height=5cm, 
	width=5cm, 
	axis lines=middle,
	minor x tick num={1},
	minor y tick num={1},
	grid=both,
	%xtick={0,50,...,250},
	%ytick={0,10,20,...,40},
	%axis line style={-latex}, 
	xmin=\xmin,xmax=\xmax, 
	ymin=\ymin,ymax=\ymax,
	%restrict y to domain=-1:10,
	%no markers, 
	xlabel={$$},ylabel={$$}, 
	%every axis x label/.append style={anchor=north}, 
	%every axis y label/.append style={anchor=east}
	]
	
	\draw[red,thick,->] (axis cs:\xa,\xb) -- (axis cs:\ya,\yb);
	\draw[blue,thick,->] (axis cs:0,0) -- (axis cs:2,3) node [right] {x};
	\draw[blue,thick,->] (axis cs:0,0) -- (axis cs:-1,1) node [left] {y};
	%\draw [fill=blue] (A) circle (2pt) node [left] {A}; 
	%\draw [fill=blue] (B) circle (2pt) node [left] {B};

	%\draw [-latex, red, thick, shorten >= 1.00cm] (A) -- (B);
	%%\addplot[color=red,solid,mark=*] coordinates {(-3,-4)};
	%%\addplot[domain=\xmin:\xmax,thick,samples=100,draw=red,<->]{-(x-2)^2+5} [yshift=8pt] node[pos=0.7,right] {$g(x)$};
	
\end{axis}
\end{tikzpicture}
\item $\mathbf{x-y}=(3,2)$\\
\renewcommand{\xa}{-1}
\renewcommand{\xb}{1}
\renewcommand{\ya}{2}
\renewcommand{\yb}{3}\renewcommand{\xmax}{5}
\renewcommand{\xmin}{-5}
\renewcommand{\xstep}{0} %must be xmin+n
\renewcommand{\ymax}{5}
\renewcommand{\ymin}{-5}
\renewcommand{\ystep}{0} %must be ymin+n
%%%%%%%
\begin{tikzpicture} 
\begin{axis}[height=5cm, 
	width=5cm, 
	axis lines=middle,
	minor x tick num={1},
	minor y tick num={1},
	grid=both,
	%xtick={0,50,...,250},
	%ytick={0,10,20,...,40},
	%axis line style={-latex}, 
	xmin=\xmin,xmax=\xmax, 
	ymin=\ymin,ymax=\ymax,
	%restrict y to domain=-1:10,
	%no markers, 
	xlabel={$$},ylabel={$$}, 
	%every axis x label/.append style={anchor=north}, 
	%every axis y label/.append style={anchor=east}
	]
	
	\draw[red,thick,->] (axis cs:\xa,\xb) -- (axis cs:\ya,\yb);
	\draw[blue,thick,->] (axis cs:0,0) -- (axis cs:2,3) node [right] {x};
	\draw[blue,thick,->] (axis cs:0,0) -- (axis cs:-1,1) node [left] {y};
	%\draw [fill=blue] (A) circle (2pt) node [left] {A}; 
	%\draw [fill=blue] (B) circle (2pt) node [left] {B};

	%\draw [-latex, red, thick, shorten >= 1.00cm] (A) -- (B);
	%%\addplot[color=red,solid,mark=*] coordinates {(-3,-4)};
	%%\addplot[domain=\xmin:\xmax,thick,samples=100,draw=red,<->]{-(x-2)^2+5} [yshift=8pt] node[pos=0.7,right] {$g(x)$};
	
\end{axis}
\end{tikzpicture}
\item $\mathbf{x+}2\mathbf{y}=(2,3)+2(-1,1)=(2-2,3+2)=(0,5)$\\
\renewcommand{\xa}{0}
\renewcommand{\xb}{0}
\renewcommand{\ya}{0}
\renewcommand{\yb}{5}\renewcommand{\xmax}{5}
\renewcommand{\xmin}{-5}
\renewcommand{\xstep}{0} %must be xmin+n
\renewcommand{\ymax}{5}
\renewcommand{\ymin}{-5}
\renewcommand{\ystep}{0} %must be ymin+n
%%%%%%%
\begin{tikzpicture} 
\begin{axis}[height=5cm, 
	width=5cm, 
	axis lines=middle,
	minor x tick num={1},
	minor y tick num={1},
	grid=both,
	%xtick={0,50,...,250},
	%ytick={0,10,20,...,40},
	%axis line style={-latex}, 
	xmin=\xmin,xmax=\xmax, 
	ymin=\ymin,ymax=\ymax,
	%restrict y to domain=-1:10,
	%no markers, 
	xlabel={$$},ylabel={$$}, 
	%every axis x label/.append style={anchor=north}, 
	%every axis y label/.append style={anchor=east}
	]
	
	\draw[red,thick,->] (axis cs:\xa,\xb) -- (axis cs:\ya,\yb);
	\draw[blue,thick,->] (axis cs:0,0) -- (axis cs:2,3) node [right] {x};
	\draw[blue,thick,->] (axis cs:0,0) -- (axis cs:-1,1) node [left] {y};
	%\draw [fill=blue] (A) circle (2pt) node [left] {A}; 
	%\draw [fill=blue] (B) circle (2pt) node [left] {B};

	%\draw [-latex, red, thick, shorten >= 1.00cm] (A) -- (B);
	%%\addplot[color=red,solid,mark=*] coordinates {(-3,-4)};
	%%\addplot[domain=\xmin:\xmax,thick,samples=100,draw=red,<->]{-(x-2)^2+5} [yshift=8pt] node[pos=0.7,right] {$g(x)$};
	
\end{axis}
\end{tikzpicture}\columnbreak
\item $\nicefrac{1}{2}\mathbf{x+}\nicefrac{1}{2}\mathbf{y}=\nicefrac{1}{2}(\mathbf{x+y})=(\nicefrac{1}{2},2)$\\
\renewcommand{\ya}{.5}
\renewcommand{\yb}{2}\renewcommand{\xmax}{5}
\renewcommand{\xmin}{-5}
\renewcommand{\xstep}{0} %must be xmin+n
\renewcommand{\ymax}{5}
\renewcommand{\ymin}{-5}
\renewcommand{\ystep}{0} %must be ymin+n
%%%%%%%
\begin{tikzpicture} 
\begin{axis}[height=5cm, 
	width=5cm, 
	axis lines=middle,
	minor x tick num={1},
	minor y tick num={1},
	grid=both,
	%xtick={0,50,...,250},
	%ytick={0,10,20,...,40},
	%axis line style={-latex}, 
	xmin=\xmin,xmax=\xmax, 
	ymin=\ymin,ymax=\ymax,
	%restrict y to domain=-1:10,
	%no markers, 
	xlabel={$$},ylabel={$$}, 
	%every axis x label/.append style={anchor=north}, 
	%every axis y label/.append style={anchor=east}
	]
	
	\draw[red,thick,->] (axis cs:\xa,\xb) -- (axis cs:\ya,\yb);
	\draw[blue,thick,->] (axis cs:0,0) -- (axis cs:2,3) node [right] {x};
	\draw[blue,thick,->] (axis cs:0,0) -- (axis cs:-1,1) node [left] {y};
	%\draw [fill=blue] (A) circle (2pt) node [left] {A}; 
	%\draw [fill=blue] (B) circle (2pt) node [left] {B};

	%\draw [-latex, red, thick, shorten >= 1.00cm] (A) -- (B);
	%%\addplot[color=red,solid,mark=*] coordinates {(-3,-4)};
	%%\addplot[domain=\xmin:\xmax,thick,samples=100,draw=red,<->]{-(x-2)^2+5} [yshift=8pt] node[pos=0.7,right] {$g(x)$};
	
\end{axis}
\end{tikzpicture}
\item $\mathbf{y-x}=-\left(\mathbf{x-y}\right)=(-3,-2)$\\
\renewcommand{\xa}{2}
\renewcommand{\xb}{3}
\renewcommand{\ya}{-1}
\renewcommand{\yb}{1}\renewcommand{\xmax}{5}
\renewcommand{\xmin}{-5}
\renewcommand{\xstep}{0} %must be xmin+n
\renewcommand{\ymax}{5}
\renewcommand{\ymin}{-5}
\renewcommand{\ystep}{0} %must be ymin+n
%%%%%%%
\begin{tikzpicture} 
\begin{axis}[height=5cm, 
	width=5cm, 
	axis lines=middle,
	minor x tick num={1},
	minor y tick num={1},
	grid=both,
	%xtick={0,50,...,250},
	%ytick={0,10,20,...,40},
	%axis line style={-latex}, 
	xmin=\xmin,xmax=\xmax, 
	ymin=\ymin,ymax=\ymax,
	%restrict y to domain=-1:10,
	%no markers, 
	xlabel={$$},ylabel={$$}, 
	%every axis x label/.append style={anchor=north}, 
	%every axis y label/.append style={anchor=east}
	]
	
	\draw[red,thick,->] (axis cs:\xa,\xb) -- (axis cs:\ya,\yb);
	\draw[blue,thick,->] (axis cs:0,0) -- (axis cs:2,3) node [right] {x};
	\draw[blue,thick,->] (axis cs:0,0) -- (axis cs:-1,1) node [left] {y};
	%\draw [fill=blue] (A) circle (2pt) node [left] {A}; 
	%\draw [fill=blue] (B) circle (2pt) node [left] {B};

	%\draw [-latex, red, thick, shorten >= 1.00cm] (A) -- (B);
	%%\addplot[color=red,solid,mark=*] coordinates {(-3,-4)};
	%%\addplot[domain=\xmin:\xmax,thick,samples=100,draw=red,<->]{-(x-2)^2+5} [yshift=8pt] node[pos=0.7,right] {$g(x)$};
	
\end{axis}
\end{tikzpicture}
\item $2\mathbf{x-y}=2(2,3)-(-1,1)=(4+1,6-1)=(5,5)$\\
\renewcommand{\xa}{0}
\renewcommand{\xb}{0}
\renewcommand{\ya}{5}
\renewcommand{\yb}{5}\renewcommand{\xmax}{5}
\renewcommand{\xmin}{-5}
\renewcommand{\xstep}{0} %must be xmin+n
\renewcommand{\ymax}{5}
\renewcommand{\ymin}{-5}
\renewcommand{\ystep}{0} %must be ymin+n
%%%%%%%
\begin{tikzpicture} 
\begin{axis}[height=5cm, 
	width=5cm, 
	axis lines=middle,
	minor x tick num={1},
	minor y tick num={1},
	grid=both,
	%xtick={0,50,...,250},
	%ytick={0,10,20,...,40},
	%axis line style={-latex}, 
	xmin=\xmin,xmax=\xmax, 
	ymin=\ymin,ymax=\ymax,
	%restrict y to domain=-1:10,
	%no markers, 
	xlabel={$$},ylabel={$$}, 
	%every axis x label/.append style={anchor=north}, 
	%every axis y label/.append style={anchor=east}
	]
	
	\draw[red,thick,->] (axis cs:\xa,\xb) -- (axis cs:\ya,\yb);
	\draw[blue,thick,->] (axis cs:0,0) -- (axis cs:2,3) node [right] {x};
	\draw[blue,thick,->] (axis cs:0,0) -- (axis cs:-1,1) node [left] {y};
	%\draw [fill=blue] (A) circle (2pt) node [left] {A}; 
	%\draw [fill=blue] (B) circle (2pt) node [left] {B};

	%\draw [-latex, red, thick, shorten >= 1.00cm] (A) -- (B);
	%%\addplot[color=red,solid,mark=*] coordinates {(-3,-4)};
	%%\addplot[domain=\xmin:\xmax,thick,samples=100,draw=red,<->]{-(x-2)^2+5} [yshift=8pt] node[pos=0.7,right] {$g(x)$};
	
\end{axis}
\end{tikzpicture}\columnbreak
\item $\left\Vert \mathbf{x}\right\Vert =\sqrt{2^{2}+3^{2}}=\sqrt{13}$\\
\renewcommand{\ya}{3.6055}
\renewcommand{\yb}{0}\renewcommand{\xmax}{5}
\renewcommand{\xmin}{-5}
\renewcommand{\xstep}{0} %must be xmin+n
\renewcommand{\ymax}{5}
\renewcommand{\ymin}{-5}
\renewcommand{\ystep}{0} %must be ymin+n
%%%%%%%
\begin{tikzpicture} 
\begin{axis}[height=5cm, 
	width=5cm, 
	axis lines=middle,
	minor x tick num={1},
	minor y tick num={1},
	grid=both,
	%xtick={0,50,...,250},
	%ytick={0,10,20,...,40},
	%axis line style={-latex}, 
	xmin=\xmin,xmax=\xmax, 
	ymin=\ymin,ymax=\ymax,
	%restrict y to domain=-1:10,
	%no markers, 
	xlabel={$$},ylabel={$$}, 
	%every axis x label/.append style={anchor=north}, 
	%every axis y label/.append style={anchor=east}
	]
	
	\draw[red,thick,->] (axis cs:\xa,\xb) -- (axis cs:\ya,\yb);
	\draw[blue,thick,->] (axis cs:0,0) -- (axis cs:2,3) node [right] {x};
	\draw[blue,thick,->] (axis cs:0,0) -- (axis cs:-1,1) node [left] {y};
	%\draw [fill=blue] (A) circle (2pt) node [left] {A}; 
	%\draw [fill=blue] (B) circle (2pt) node [left] {B};

	%\draw [-latex, red, thick, shorten >= 1.00cm] (A) -- (B);
	%%\addplot[color=red,solid,mark=*] coordinates {(-3,-4)};
	%%\addplot[domain=\xmin:\xmax,thick,samples=100,draw=red,<->]{-(x-2)^2+5} [yshift=8pt] node[pos=0.7,right] {$g(x)$};
	
\end{axis}
\end{tikzpicture}
\item $\nicefrac{\mathbf{X}}{\Vert\mathbf{x}\Vert}=\nicefrac{1}{\sqrt{13}}(2,3)=(\nicefrac{2}{\sqrt{13}},\nicefrac{3}{\sqrt{13}})$\\
\renewcommand{\ya}{.5547}
\renewcommand{\yb}{.83205}\renewcommand{\xmax}{5}
\renewcommand{\xmin}{-5}
\renewcommand{\xstep}{0} %must be xmin+n
\renewcommand{\ymax}{5}
\renewcommand{\ymin}{-5}
\renewcommand{\ystep}{0} %must be ymin+n
%%%%%%%
\begin{tikzpicture} 
\begin{axis}[height=5cm, 
	width=5cm, 
	axis lines=middle,
	minor x tick num={1},
	minor y tick num={1},
	grid=both,
	%xtick={0,50,...,250},
	%ytick={0,10,20,...,40},
	%axis line style={-latex}, 
	xmin=\xmin,xmax=\xmax, 
	ymin=\ymin,ymax=\ymax,
	%restrict y to domain=-1:10,
	%no markers, 
	xlabel={$$},ylabel={$$}, 
	%every axis x label/.append style={anchor=north}, 
	%every axis y label/.append style={anchor=east}
	]
	
	\draw[blue,thick,->] (axis cs:0,0) -- (axis cs:2,3) node [right] {x};
	\draw[blue,thick,->] (axis cs:0,0) -- (axis cs:-1,1) node [left] {y};
	\draw[red,thick,->] (axis cs:\xa,\xb) -- (axis cs:\ya,\yb);
	%\draw [fill=blue] (A) circle (2pt) node [left] {A}; 
	%\draw [fill=blue] (B) circle (2pt) node [left] {B};

	%\draw [-latex, red, thick, shorten >= 1.00cm] (A) -- (B);
	%%\addplot[color=red,solid,mark=*] coordinates {(-3,-4)};
	%%\addplot[domain=\xmin:\xmax,thick,samples=100,draw=red,<->]{-(x-2)^2+5} [yshift=8pt] node[pos=0.7,right] {$g(x)$};
	
\end{axis}
\end{tikzpicture}\\
This is the \emph{unit vector} as $\Vert\nicefrac{\mathbf{X}}{\Vert\mathbf{x}\Vert}\Vert=1$.
\end{enumerate}
\end{multicols}
\item [\textbf{5.}]If a point $(x_{1},x_{2})$ is on the line $\ell$ then
$x_{1}=1-2t$ and $x_{2}=3+t$ for some $t\in\mathbb{R}$.

\begin{enumerate}[resume]
\item $x_{1}=-1\Rightarrow t=1$ so $x_{2}=3+1=4$. This point lies on $\ell$.
\item $x_{2}=0\Rightarrow t=-3$ so $x_{1}=1-2(-3)=7$. This point lies
on $\ell$.
\item $x_{2}=2\Rightarrow t=-1$ so $x_{1}=3$. This point does not lie
on $\ell$.
\end{enumerate}
\item [\textbf{6.}](There may be multiple solutions)

\begin{enumerate}
\item If $3x_{1}+4x_{2}=6$ then $x_{1}=-\nicefrac{4}{3}x_{2}+2$. Thus
$\mathbf{x}=(x_{1},x_{2})=(-\nicefrac{4}{3}x_{2}+2,x_{2})=x_{2}(-\nicefrac{4}{3},1)+(2,0)$.
Replacing $x_{2}=t$, $\mathbf{x=}t(-\nicefrac{4}{3},1)+(2,0)$
\item The slope $\nicefrac{1}{3}$ gives the direction vector $(3,1)$ for
a line that starts at $(-1,2)$, $\mathbf{x}=(-1,2)+t(3,1)$.
\item [\textbf{d.}]The direction vector is $(3,5)$ for a line that starts
at $(-2,1)$, $\mathbf{x}=(-2,1)+t(3,5)$. 
\end{enumerate}
\item [\textbf{10.a.}] The plane containing the point $(-1,0,1)$ and the
line $\mathbf{x}=(1,1,1)+t(1,7,-1)$ also contains the point $(1,1,1)$
(just let $t=0$) and is parallel to the vector $(1,7,-1)$. 

\begin{itemize}
\item We need two non-parallel lines which are both parallel to the plane.
To find another vector note that the difference between the vectors,
$(1,1,1)-(-1,0,1)=(2,1,0)$, must be parallel to the plane since both
these points are on the plane. 
\end{itemize}
Now $\mathbf{x}=(1,1,1)+t(1,7,-1)+s(2,1,0)$, for $s,t\in\mathbb{R}$.
\item [\textbf{13.}]$\overrightarrow{AM}=\frac{1}{2}\overrightarrow{AB}$
and $\overrightarrow{AN}=\frac{1}{2}\overrightarrow{AC}$. 

So then $\overrightarrow{MN}=\overrightarrow{AN}-\overrightarrow{AM}=\frac{1}{2}\overrightarrow{AC}-\frac{1}{2}\overrightarrow{AB}=\frac{1}{2}\left(\overrightarrow{AC}-\overrightarrow{AB}\right)=\frac{1}{2}\overrightarrow{BC}$.
\item [\textbf{22.}]

\begin{enumerate}
\item $\mathbf{v}=c_{1}\mathbf{v_{1}}+c_{2}\mathbf{v_{2}}+\dots+c_{2}\mathbf{v_{k}}$
where $c_{i}\in\mathbb{R}$ for each $i=1,\dots,k$. So then for any
$c\in\mathbb{R}$, $c\mathbf{v}=cc_{1}\mathbf{v_{1}}+cc_{2}\mathbf{v_{2}}+\dots+cc_{k}\mathbf{v_{k}}$
is a linear combination of $\mathbf{v}_{1},\dots,\mathbf{v}_{k}$
since $cc_{i}\in\mathbb{R}$ for each $i=1,\dots,k$.
\item $\mathbf{v}=c_{1}\mathbf{v_{1}}+c_{2}\mathbf{v_{2}}+\dots+c_{2}\mathbf{v_{k}}$
and $\mathbf{w}=d_{1}\mathbf{v_{1}}+d_{2}\mathbf{v_{2}}+\dots+d_{k}\mathbf{v_{k}}$
where $c_{i},d_{i}\in\mathbb{R}$ for each $i=1,\dots,k$. So then
\begin{eqnarray*}
\mathbf{v+w} & = & c_{1}\mathbf{v_{1}}+c_{2}\mathbf{v_{2}}+\dots+c_{k}\mathbf{v_{k}}+d_{1}\mathbf{v_{1}}+d_{2}\mathbf{v_{2}}+\dots+d_{k}\mathbf{v_{k}}\\
 & = & c_{1}\mathbf{v_{1}}+d_{1}\mathbf{v_{1}}+c_{2}\mathbf{v_{2}}+d_{2}\mathbf{v_{2}}+\dots+d_{k}\mathbf{v_{k}}\dots+d_{k}\mathbf{v_{k}}\\
 & = & (c_{1}+d_{1})\mathbf{v_{1}}+(c_{2}+d_{2})\mathbf{v_{2}}+\dots+(c_{k}+d_{k})\mathbf{v_{k}}
\end{eqnarray*}
Which is a linear combination of $\mathbf{v}_{1},\dots,\mathbf{v}_{k}$. 
\end{enumerate}
\end{enumerate}

\end{document}
